% Options for packages loaded elsewhere
\PassOptionsToPackage{unicode}{hyperref}
\PassOptionsToPackage{hyphens}{url}
\documentclass[
]{article}
\usepackage{ma2003b-notes}
\hypersetup{
  pdftitle={Discriminant Analysis - Companion Notes},
  pdfauthor={Dr. Juliho Castillo},
  hidelinks,
  pdfcreator={LaTeX via pandoc}}

\title{Discriminant Analysis - Companion Notes}
\author{Dr. Juliho Castillo}
\date{}

\begin{document}
\maketitle

\textbf{Discriminant Analysis}

Companion Notes for Chapter 5

MA2003B - Multivariate Methods in Data Science

Dr. Juliho Castillo ⋅ Tecnologico de Monterrey

\begin{center}\rule{0.5\linewidth}{0.5pt}\end{center}

\textbf{Hands-On Learning Resource}

This document is accompanied by an interactive Jupyter notebook that
demonstrates all concepts with executable Python code and
visualizations:

\texttt{ch5\_guiding\_example/marketing\_discriminant\_analysis.ipynb}

The notebook analyzes a customer segmentation case study with 1,200
e-commerce customers, comparing Linear Discriminant Analysis (LDA) and
Quadratic Discriminant Analysis (QDA) for marketing campaign
optimization.

Each major concept in these notes references specific modules in the
notebook where you can see the implementation and results.

\section{Introduction: The Classification Problem}

\subsection{Why Discriminant Analysis?}

Imagine you are a marketing analyst at an e-commerce company with
thousands of customers. You want to answer a fundamental question:
\textbf{Given a customer's behavior, which segment do they belong to?}
Are they a high-value customer worth premium retention efforts? A loyal
regular who might respond to upselling? Or an occasional buyer who needs
re-engagement?

This is a \textbf{classification problem} - we have several known groups
and want to assign new observations to the correct group based on
measured characteristics. Discriminant Analysis provides a principled
statistical approach to this problem.

\subsection{Real-World Applications}

Discriminant analysis finds extensive application across diverse domains
where classification decisions require both accuracy and
interpretability. In business and marketing contexts, the technique
enables customer segmentation for targeted campaigns by identifying
behavioral patterns that distinguish high-value customers from
occasional buyers. Financial institutions employ discriminant analysis
for credit risk assessment, where the method helps determine whether to
approve or reject loan applications based on applicant characteristics
and historical default patterns. Additionally, churn prediction models
leverage discriminant functions to identify customers at risk of
discontinuing service, allowing proactive retention efforts.

The healthcare and medical sectors utilize discriminant analysis for
diagnostic purposes, where the method combines patient symptoms,
laboratory test results, and medical history to classify individuals
into disease categories. Treatment response prediction represents
another critical application, as clinicians can use discriminant models
to identify which patients are likely to benefit from specific
therapeutic interventions. Medical imaging classification similarly
benefits from discriminant techniques, particularly when distinguishing
between normal and pathological tissue patterns in radiological studies.

Manufacturing industries apply discriminant analysis extensively in
quality control processes, where products must be classified as
acceptable, borderline, or defective based on multiple measurement
criteria. Defect type classification helps identify the root causes of
manufacturing problems by associating specific defects with process
parameters. Process monitoring and fault detection systems employ
discriminant functions to distinguish normal operating conditions from
various failure modes, enabling timely intervention before quality
degradation occurs.

In sports and performance contexts, discriminant analysis assists in
athlete classification for training programs by identifying
physiological and performance characteristics that predict success in
different athletic domains. Talent identification systems use
discriminant models to evaluate prospective athletes based on multiple
attributes, helping coaches allocate resources effectively. Performance
level assessment employs these techniques to classify athletes into
skill categories, informing coaching strategies and competitive
placement decisions.

\subsection{The Core Idea}

Discriminant Analysis finds \textbf{discriminant functions} --- linear
(or quadratic) combinations of your predictor variables that best
separate your groups. Think of it as finding the ``best viewing angle''
to distinguish between groups in multidimensional space.

\textbf{Key Question}: Which combination of behavioral metrics (purchase
frequency, order value, engagement, etc.) best distinguishes high-value
customers from loyal customers from occasional customers?

\section{Mathematical Foundations}

\subsection{The Setup}

The discriminant analysis framework requires formalization of the
classification problem structure. The analysis involves \(g\) distinct
groups or populations into which observations must be classified, with
\(p\) predictor variables measured on each observation. A training
dataset provides examples of observations with known group memberships,
enabling the estimation of discriminant functions. The ultimate goal
consists of developing a classification rule that assigns new
observations to one of the \(g\) groups based on their predictor values.

The mathematical notation employed throughout discriminant analysis
includes several key quantities. The vector
\(\mathbf{x} = \left( x_{1},x_{2},\ldots,x_{p} \right)^{\top}\)
represents the predictor variables for a given observation, where the
superscript indicates vector transposition. The prior probability
\(\pi_{k}\) denotes the probability that a randomly selected observation
belongs to group \(k\), which may reflect either the proportion of group
\(k\) in the population or be specified based on domain considerations.
The mean vector \(\mathbf{\mu}_{k}\) characterizes the central tendency
of group \(k\) in the \(p\)-dimensional predictor space, while the
covariance matrix \(\mathbf{\Sigma}_{k}\) describes the variability and
correlation structure within group \(k\). Finally, the probability
density function \(f_{k}\left( \mathbf{x} \right)\) specifies the
likelihood of observing predictor values \(\mathbf{x}\) given membership
in group \(k\), with the multivariate normal distribution serving as the
standard assumption in classical discriminant analysis.

\subsection{Classification Rules: Bayes Theorem}

The foundation of discriminant analysis is Bayes theorem, which provides
the optimal classification framework when distributional assumptions
hold. Understanding this theorem in depth clarifies how discriminant
analysis achieves its classification decisions and why these decisions
are optimal under certain conditions.

\subsubsection{Understanding the Notation}

Before examining the theorem itself, we must clarify the mathematical
notation employed throughout discriminant analysis:

\textbf{Group Membership Variable G}: The random variable \(G\)
represents the group or class to which an observation belongs. This
categorical variable can take values \(1,2,\ldots,g\), where \(g\)
denotes the total number of groups in the classification problem. For
example, in a credit risk application, \(G\) might take values from the
set \{0, 1\}, where 0 indicates ``no default'' and 1 indicates
``default''. In the marketing segmentation example, \(G\) takes values
from \{1, 2, 3\}, corresponding to High-Value, Loyal, and Occasional
customer segments respectively.

\textbf{Predictor Vector x}: The vector
\(\mathbf{x} = \left( x_{1},x_{2},\ldots,x_{p} \right)^{\top}\) contains
the values of \(p\) predictor variables measured on a specific
observation. In the marketing context, this might be
\(\mathbf{x} = \left( \text{purchase frequency},\text{ order value},\text{ browsing time},\ldots \right)^{\top}\).
Each component \(x_{j}\) represents a measured characteristic that
potentially helps discriminate between groups.

\textbf{Prior Probability pi\_k}: The quantity \(\pi_{k}\) denotes the
prior probability that a randomly selected observation belongs to group
\(k\), calculated before observing the predictor values. This
probability reflects either population frequencies or domain knowledge
about group prevalence. The prior probabilities must satisfy
\(\pi_{k} > 0\) for all groups and \(\sum_{k = 1}^{g}\pi_{k} = 1\).

\textbf{Class-Conditional Density f\_k(x)}: The function
\(f_{k}\left( \mathbf{x} \right)\) represents the probability density of
observing predictor values \(\mathbf{x}\) given that the observation
belongs to group \(k\). In discriminant analysis, we typically assume
\(f_{k}\left( \mathbf{x} \right)\) follows a multivariate normal
distribution with mean vector \(\mathbf{\mu}_{k}\) and covariance matrix
\(\mathbf{\Sigma}_{k}\). This density quantifies how typical or likely
the observed predictor values are for members of group \(k\).

\subsubsection{The Bayes Theorem Formula}

Given observation \(\mathbf{x}\), the posterior probability of belonging
to group \(k\) is:

\[P\left( G = k~|~\mathbf{x} \right) = \frac{f_{k}\left( \mathbf{x} \right)\pi_{k}}{\sum_{j = 1}^{g}f_{j}\left( \mathbf{x} \right)\pi_{j}}\]

This equation decomposes into several interpretable components:

\textbf{Left Side - Posterior Probability}: The quantity
\(P\left( G = k~|~\mathbf{x} \right)\) represents the probability that
an observation with observed predictor values \(\mathbf{x}\) belongs to
group \(k\). This is precisely what we want to calculate for
classification purposes. The vertical bar notation \(|\) reads as
``given'' or ``conditional on'', so we interpret this as ``the
probability of group membership \(k\) given that we have observed
predictor values \(\mathbf{x}\)''.

\textbf{Numerator Components}:

\begin{itemize}
\item
  \(f_{k}\left( \mathbf{x} \right)\): The likelihood of observing these
  predictor values if the observation belongs to group \(k\)
\item
  \(\pi_{k}\): The prior probability of group \(k\) before observing any
  predictor values
\item
  Product \(f_{k}\left( \mathbf{x} \right)\pi_{k}\): The joint
  probability of both belonging to group \(k\) and having predictor
  values \(\mathbf{x}\)
\end{itemize}

\textbf{Denominator}: The sum
\(\sum_{j = 1}^{g}f_{j}\left( \mathbf{x} \right)\pi_{j}\) aggregates the
joint probabilities across all \(g\) groups. This normalization constant
ensures that the posterior probabilities sum to unity across all
possible group assignments:
\(\sum_{k = 1}^{g}P\left( G = k~|~\mathbf{x} \right) = 1\).

\subsubsection{Intuitive Interpretation}

Bayes theorem combines two distinct sources of information to produce
classification decisions:

\textbf{Prior Knowledge}: The prior probabilities \(\pi_{k}\) encode
what we know about group frequencies before examining any specific
observation. If historical data shows that 5 percent of loan applicants
default, we set \(\pi_{1} = 0.05\) for the default group. This prior
information prevents the classifier from ignoring base rates and making
unrealistic predictions.

\textbf{Evidence from Data}: The class-conditional density
\(f_{k}\left( \mathbf{x} \right)\) quantifies how typical the observed
predictor values are for group \(k\). If a loan applicant has income
equal to 50,000 dollars, debt ratio equal to 0.4, and credit score equal
to 650, the density \(f_{1}\left( \mathbf{x} \right)\) tells us how
likely these values are among customers who ultimately default, while
\(f_{0}\left( \mathbf{x} \right)\) tells us how likely they are among
customers who do not default.

The posterior probability is proportional to the product of these two
components: ``how likely this data is for group \(k\)'' multiplied by
``how common group \(k\) is''. Groups that are both common (high
\(\pi_{k}\)) and compatible with the observed data (high
\(f_{k}\left( \mathbf{x} \right)\)) receive high posterior
probabilities.

\subsubsection{Detailed Example: Credit Risk Assessment}

Consider a credit risk classification problem with the following
characteristics:

\textbf{Groups}: \(k = 0\) (no default), \(k = 1\) (default)

\textbf{Prior Probabilities}: Based on historical data,
\(\pi_{0} = 0.95\) (95 percent of customers do not default) and
\(\pi_{1} = 0.05\) (5 percent default)

\textbf{Observation}: A loan applicant with
\(\mathbf{x} = \left\lbrack \text{income } = 50000,\text{ debt ratio } = 0.4,\text{ credit score } = 650 \right\rbrack^{\top}\)

\textbf{Class-Conditional Densities}: Suppose our fitted discriminant
analysis model produces:

\begin{itemize}
\item
  \(f_{0}\left( \mathbf{x} \right) = 0.0008\) (this financial profile is
  somewhat typical for non-defaulters)
\item
  \(f_{1}\left( \mathbf{x} \right) = 0.0030\) (this financial profile is
  more typical for defaulters)
\end{itemize}

The posterior probabilities calculate as:

Numerator for group 0:
\(f_{0}\left( \mathbf{x} \right)\pi_{0} = 0.0008 \times 0.95 = 0.00076\)

Numerator for group 1:
\(f_{1}\left( \mathbf{x} \right)\pi_{1} = 0.0030 \times 0.05 = 0.00015\)

Denominator:
\(\sum_{j = 0}^{1}f_{j}\left( \mathbf{x} \right)\pi_{j} = 0.00076 + 0.00015 = 0.00091\)

Posterior probabilities:

\begin{itemize}
\item
  \(P\left( G = 0~|~\mathbf{x} \right) = \frac{0.00076}{0.00091} \approx 0.835\)
  (83.5 percent probability of no default)
\item
  \(P\left( G = 1~|~\mathbf{x} \right) = \frac{0.00015}{0.00091} \approx 0.165\)
  (16.5 percent probability of default)
\end{itemize}

Despite the financial profile being more typical of defaulters (density
ratio \(\frac{f_{1}}{f_{0}} = 3.75\)), the strong prior probability
favoring non-default prevents the classifier from predicting default.
The model balances the evidence from the data against the base rate
information.

\subsubsection{The Bayes Classification Rule}

\textbf{Bayes Classification Rule}: Assign observation \(\mathbf{x}\) to
the group \(k^{\ast}\) that maximizes the posterior probability:

\[k^{\ast} = \arg\max\limits_{k}P\left( G = k~|~\mathbf{x} \right) = \arg\max\limits_{k}f_{k}\left( \mathbf{x} \right)\pi_{k}\]

The second equality follows because the denominator
\(\sum_{j = 1}^{g}f_{j}\left( \mathbf{x} \right)\pi_{j}\) is identical
across all groups and therefore does not affect which group achieves the
maximum. Consequently, we need only compare the numerators
\(f_{k}\left( \mathbf{x} \right)\pi_{k}\) across groups.

In the credit risk example above, we would classify the applicant as
non-default (group 0) because
\(P\left( G = 0~|~\mathbf{x} \right) = 0.835 > P\left( G = 1~|~\mathbf{x} \right) = 0.165\).

\subsubsection{Optimality of the Bayes Rule}

The Bayes classification rule is optimal in a precise mathematical
sense: it minimizes the total probability of misclassification. This
result, known as the Bayes risk minimization theorem, states that among
all possible classification rules, assigning observations to the group
with the highest posterior probability achieves the lowest possible
error rate.

This optimality holds under three critical conditions:

\textbf{Condition 1 - Correct Prior Probabilities}: The specified prior
probabilities \(\pi_{k}\) accurately reflect the true prevalence of
groups in the population or appropriately weight the relative importance
of different groups.

\textbf{Condition 2 - Correct Density Specification}: The assumed
class-conditional densities \(f_{k}\left( \mathbf{x} \right)\) correctly
specify the distribution of predictor values within each group. In
discriminant analysis, we assume multivariate normal distributions, so
optimality requires that this assumption holds or approximately holds.

\textbf{Condition 3 - Equal Misclassification Costs (0-1 Loss)}:
``Minimizes the total probability of misclassification'' and ``assign to
the highest posterior probability'' are only the same rule under 0-1
loss, where every type of error counts equally. Under asymmetric costs
(e.g. missing a default is worse than a false alarm), the
\textbf{risk}-minimizing rule instead compares
\(\pi_{k}f_{k}\left( \mathbf{x} \right) \cdot C_{k}\) weighted by cost,
and can assign an observation to a group whose posterior probability is
not the largest.

When these conditions are violated, the Bayes rule using the
misspecified priors or densities may perform suboptimally. Nevertheless,
the framework remains valuable because:

\textbf{Robustness}: Discriminant analysis often performs well even when
the normality assumption is violated moderately, particularly with large
sample sizes.

\textbf{Flexibility in Priors}: We can adjust prior probabilities to
reflect business costs or domain knowledge, even if these differ from
sample proportions.

\textbf{Interpretability}: The probabilistic framework provides
posterior probabilities that quantify classification uncertainty,
enabling risk-based decision making beyond simple classification.

\subsubsection{Connection to Discriminant Functions}

The discriminant analysis methods (LDA and QDA) implement the Bayes
classification rule by making specific assumptions about the
class-conditional densities \(f_{k}\left( \mathbf{x} \right)\). Both
methods assume multivariate normal distributions:

\[f_{k}\left( \mathbf{x} \right) = \frac{1}{(2\pi)^{p/2}|\mathbf{\Sigma}_{k}|^{1/2}}\exp( - \frac{1}{2}\left( \mathbf{x} - \mathbf{\mu}_{k} \right)^{\top}\mathbf{\Sigma}_{k}^{- 1}\left( \mathbf{x} - \mathbf{\mu}_{k} \right))\]

Taking logarithms and simplifying yields discriminant functions, which
are equivalent to comparing posterior probabilities:

\textbf{LDA}: Assuming \(\mathbf{\Sigma}_{k} = \mathbf{\Sigma}\) for all
groups yields linear discriminant functions

\textbf{QDA}: Allowing group-specific \(\mathbf{\Sigma}_{k}\) yields
quadratic discriminant functions

Thus, discriminant analysis provides a computationally efficient
implementation of the optimal Bayes rule under normality assumptions,
transforming the classification problem into comparison of discriminant
scores rather than explicit probability calculations.

\subsection{Linear Discriminant Analysis (LDA)}

\subsubsection{Assumptions}

LDA makes two critical assumptions:

\begin{enumerate}
\item
  \textbf{Multivariate Normality}: Each group follows a multivariate
  normal distribution
\item
  \textbf{Equal Covariances}: All groups share the same covariance
  matrix:
  \(\mathbf{\Sigma}_{1} = \mathbf{\Sigma}_{2} = \ldots = \mathbf{\Sigma}_{g} = \mathbf{\Sigma}\)
\end{enumerate}

Under these assumptions, the discriminant functions become
\textbf{linear} in \(\mathbf{x}\).

\subsubsection{Fisher Linear Discriminant}

An alternative (but equivalent) approach by R.A. Fisher: Find linear
combinations of variables that maximize the ratio of between-group
variance to within-group variance.

\textbf{For two groups}, find weights \(\mathbf{a}\) that maximize:

\[\text{ maximize }\quad\frac{\left( {\overline{y}}_{1} - {\overline{y}}_{2} \right)^{2}}{s_{1}^{2} + s_{2}^{2}}\]

where \(y = \mathbf{a}^{\top}\mathbf{x}\) is the discriminant score.

The solution is:
\(\mathbf{a} \propto \mathbf{\Sigma}^{- 1}\left( \mathbf{\mu}_{1} - \mathbf{\mu}_{2} \right)\)

This gives us the \textbf{linear discriminant function}.

\subsubsection{Discriminant Scores}

For observation \(\mathbf{x}\), the discriminant score for group \(k\)
is:

\[\delta_{k}\left( \mathbf{x} \right) = \mathbf{x}^{\top}\mathbf{\Sigma}^{- 1}\mathbf{\mu}_{k} - \frac{1}{2}\mathbf{\mu}_{k}^{\top}\mathbf{\Sigma}^{- 1}\mathbf{\mu}_{k} + \log(\pi_{k})\]

\textbf{Classification Rule}: Assign \(\mathbf{x}\) to the group with
the largest discriminant score \(\delta_{k}\left( \mathbf{x} \right)\).

\subsubsection{Geometric Interpretation}

The geometric perspective on linear discriminant analysis provides
valuable intuition about the classification mechanism. Each discriminant
function defines a hyperplane in the \(p\)-dimensional predictor space,
with these hyperplanes serving as decision boundaries that partition the
space into regions corresponding to different group assignments. The
linearity of these decision boundaries, which gives the method its name,
implies that the boundaries consist of straight lines in two dimensions,
flat planes in three dimensions, and hyperplanes in higher dimensions.
The boundary between groups \(k\) and \(l\) is the hyperplane with
normal vector
\(\mathbf{\Sigma}^{- 1}\left( \mathbf{\mu}_{k} - \mathbf{\mu}_{l} \right)\),
passing through the point equidistant (in \textbf{Mahalanobis} distance)
between the two centroids when priors are equal. This normal vector
equals \(\mathbf{\mu}_{k} - \mathbf{\mu}_{l}\) itself -- making the
boundary a literal perpendicular bisector in ordinary Euclidean
coordinates -- only in the special case where \(\mathbf{\Sigma}\) is
proportional to the identity matrix (uncorrelated predictors with equal
variances). In general, \(\mathbf{\Sigma}^{- 1}\) tilts and rescales the
boundary, so ``nearest centroid'' must be read in Mahalanobis distance,
not raw Euclidean distance.

\subsection{Quadratic Discriminant Analysis (QDA)}

\subsubsection{When LDA Is Not Enough}

QDA relaxes the equal covariance assumption. Each group \(k\) has its
own covariance matrix \(\mathbf{\Sigma}_{k}\).

\textbf{When to use QDA}:

\begin{itemize}
\item
  Groups have genuinely different variability patterns
\item
  You have sufficient sample size (more parameters to estimate)
\item
  Linear boundaries do not fit the data well
\end{itemize}

\subsubsection{Quadratic Discriminant Functions}

The discriminant score becomes:

\[\delta_{k}\left( \mathbf{x} \right) = - \frac{1}{2}\log|\mathbf{\Sigma}_{k}| - \frac{1}{2}\left( \mathbf{x} - \mathbf{\mu}_{k} \right)^{\top}\mathbf{\Sigma}_{k}^{- 1}\left( \mathbf{x} - \mathbf{\mu}_{k} \right) + \log(\pi_{k})\]

This is \textbf{quadratic} in \(\mathbf{x}\), leading to curved
(quadratic) decision boundaries.

\subsubsection{Trade-offs: LDA vs QDA}

The choice between Linear Discriminant Analysis and Quadratic
Discriminant Analysis involves balancing model complexity against
estimation precision. LDA offers several compelling advantages that make
it the preferred starting point for most applications. The method
requires estimation of only \(\frac{p(p + 1)}{2}\) covariance parameters
for the pooled covariance matrix, compared to
\(g \cdot \frac{p(p + 1)}{2}\) parameters required by QDA when each of
the \(g\) groups maintains its own covariance structure. This parameter
efficiency translates to greater stability when working with smaller
sample sizes, as fewer quantities require estimation from limited data.
The reduced complexity also makes LDA less prone to overfitting,
particularly when the number of predictors approaches the sample size.
Additionally, the linear structure of LDA discriminant functions
facilitates interpretation, as the coefficients directly indicate how
variables combine to separate groups.

Conversely, QDA provides advantages when the assumption of equal
covariances proves untenable. The method's flexibility allows it to fit
complex, curved decision boundaries that better capture the structure of
data when groups genuinely exhibit different variability patterns. This
flexibility translates to improved classification accuracy when the
equal covariance assumption is violated substantially. Furthermore, QDA
does not constrain groups to share variance structures, making it more
appropriate when domain knowledge suggests heterogeneous within-group
variability.

A practical rule of thumb guides the selection between these methods.
The recommended approach begins with LDA as the baseline model. Three
conditions suggest moving to QDA warrants consideration. First, the
sample size should be substantially larger than the number of
predictors, as QDA's additional parameters require more data for
reliable estimation. Second, exploratory analysis should reveal clearly
different spread patterns across groups, suggesting that the equal
covariance assumption may not hold. Third, LDA should demonstrate
unsatisfactory performance, with classification accuracy falling short
of acceptable thresholds or exhibiting systematic patterns of
misclassification that curved boundaries might address.

\section{Practical Implementation}

\textbf{Interactive Tutorial}: The notebook
\texttt{ch5\_guiding\_example/marketing\_discriminant\_analysis.ipynb}
demonstrates each step of this workflow with executable Python code,
detailed explanations, and visualizations.

\subsection{The Analysis Workflow}

\subsubsection{Step 1: Data Preparation}

The data preparation phase establishes the foundation for reliable
discriminant analysis, with Module 2 of the accompanying notebook
providing complete implementation details. Feature selection represents
the initial critical decision, requiring analysts to identify predictors
that effectively discriminate between groups while avoiding redundancy.
The selection process should eliminate highly correlated predictors to
mitigate multicollinearity issues that can destabilize coefficient
estimates and inflate standard errors. Domain knowledge plays an
essential role in this process, as subject matter expertise often
reveals which variables possess theoretical relevance for group
separation beyond what correlation analysis alone might suggest.

Standardization becomes necessary when variables exist on disparate
scales, as the raw magnitude differences can cause variables with larger
numerical ranges to dominate the discriminant functions inappropriately.
For instance, mixing monetary variables measured in dollars (ranging
from 0 to 200) with rate variables expressed as proportions (ranging
from 0 to 1) without standardization would give undue weight to the
monetary variables. The notebook implementation employs the
\texttt{StandardScaler} transformation to convert all features to have
mean zero and standard deviation one, ensuring that each variable
contributes to the analysis based on its discriminatory power rather
than its measurement scale.

The train-test split procedure requires careful attention to maintain
the integrity of model evaluation. Stratified sampling preserves the
proportional representation of each group in both the training and
testing subsets, preventing bias that could arise if certain groups were
over-represented in one subset. The typical split allocates seventy
percent of observations to the training set for parameter estimation and
thirty percent to the testing set for unbiased performance evaluation.
The implementation utilizes the \texttt{train\_test\_split} function
with the \texttt{stratify=y} parameter, which ensures all groups appear
in both subsets according to their original proportions.

\subsubsection{Step 2: Assumption Checking}

Verification of underlying assumptions ensures the validity of
discriminant analysis results and guides methodological choices. The
multivariate normality assumption, which posits that observations within
each group follow a multivariate normal distribution, can be assessed
through several diagnostic approaches. Q-Q plots constructed for each
variable within each group provide visual evidence of univariate
normality, while multivariate tests such as the Mardia test and
Henze-Zirkler test offer formal statistical evaluation of multivariate
normality. Fortunately, discriminant analysis demonstrates robustness to
moderate violations of normality when sample sizes are large, as the
Central Limit Theorem ensures that parameter estimates remain
approximately normally distributed even when the underlying data deviate
somewhat from normality.

The equal covariance assumption, which is fundamental to LDA, merits
particular scrutiny. Box's M test provides a formal statistical test of
covariance homogeneity across groups, though practitioners should note
that this test exhibits considerable sensitivity and frequently rejects
the null hypothesis even when differences are practically negligible.
Visual inspection of covariance matrices computed separately for each
group often provides more practical insight, allowing analysts to judge
whether observed differences justify the additional complexity of QDA.
When the equal covariance assumption is violated substantially, the
analyst faces a choice between transitioning to QDA, which accommodates
group-specific covariances, or applying variance-stabilizing
transformations to the data.

Multicollinearity assessment protects against numerical instability and
interpretation difficulties arising from highly correlated predictors.
Examination of the correlation matrix among predictors reveals pairwise
relationships, while Variance Inflation Factors quantify the degree to
which each predictor's variance is inflated due to linear relationships
with other predictors. When multicollinearity proves problematic, the
analyst should remove redundant predictors, retaining those with greater
theoretical relevance or stronger discriminatory power.

\subsubsection{Step 3: Model Fitting}

\textbf{Notebook Reference}: Module 3 (LDA) and Module 6 (QDA)
demonstrate model fitting with scikit-learn.

Linear Discriminant Analysis estimation in Python employs the
scikit-learn library's \texttt{LinearDiscriminantAnalysis} class. The
estimation procedure fits the model to training data, producing several
key outputs that inform interpretation and prediction. The
\texttt{lda.scalings\_} attribute contains the discriminant function
coefficients, which indicate how the original variables combine to form
discriminant scores. The \texttt{lda.means\_} attribute stores the group
centroids, representing the mean position of each group in the predictor
space. While the pooled covariance matrix remains internal to the
implementation, the \texttt{lda.priors\_} attribute reveals the prior
probabilities assigned to each group, and the
\texttt{lda.explained\_variance\_ratio\_} attribute quantifies how much
between-group variance each discriminant function captures.

\begin{Shaded}
\begin{Highlighting}[]
\ImportTok{from}\NormalTok{ sklearn.discriminant\_analysis }\ImportTok{import}\NormalTok{ LinearDiscriminantAnalysis}

\NormalTok{lda }\OperatorTok{=}\NormalTok{ LinearDiscriminantAnalysis()}
\NormalTok{lda.fit(X\_train, y\_train)}
\end{Highlighting}
\end{Shaded}

Quadratic Discriminant Analysis follows a parallel implementation
pattern through scikit-learn's \texttt{Quadratic\allowbreak{}Discriminant\allowbreak{}Analysis}
class. The model fitting process estimates group-specific parameters,
with the \texttt{qda.means\_} attribute providing group centroids
analogous to those in LDA. The group-specific covariance matrices, which
distinguish QDA from LDA, remain internal to the implementation but
govern the quadratic decision boundaries. The \texttt{qda.priors\_}
attribute contains prior probabilities, while the
\texttt{predict\_proba()} method generates posterior probabilities that
quantify classification confidence for each observation and group
combination.

\begin{Shaded}
\begin{Highlighting}[]
\ImportTok{from}\NormalTok{ sklearn.discriminant\_analysis }\ImportTok{import}\NormalTok{ QuadraticDiscriminantAnalysis}

\NormalTok{qda }\OperatorTok{=}\NormalTok{ QuadraticDiscriminantAnalysis()}
\NormalTok{qda.fit(X\_train, y\_train)}
\end{Highlighting}
\end{Shaded}

\subsubsection{Step 4: Interpretation}

\textbf{Notebook Reference}: Module 4 provides detailed coefficient
interpretation with group means analysis.

\textbf{Discriminant Functions}

For LDA with \(g\) groups, we get up to \(\min(g - 1,p)\) discriminant
functions. The maximum number of discriminant functions is determined by
the formula:

\[m = \min(g - 1,p)\]

where \(g\) is the number of groups and \(p\) is the number of
predictors. This limitation arises because:

\begin{itemize}
\item
  With \(g\) groups, there are only \(g - 1\) independent contrasts
  between group means
\item
  With \(p\) predictors, the discriminant space has at most \(p\)
  dimensions
\end{itemize}

For example, with 3 customer segments and 8 predictors, we obtain
\(\min(3 - 1,8) = 2\) discriminant functions.

\textbf{First discriminant function} (LD1): Explains most between-group
variation \textbf{Second discriminant function} (LD2): Second-most
variation (orthogonal to LD1)

\textbf{Marketing Example Interpretation} (from the notebook):

\begin{itemize}
\item
  \textbf{LD1} (95.8\% of variance) separates High-Value from Occasional
  customers

  \begin{itemize}
  \item
    Strong negative coefficients: purchase frequency, loyalty points
  \item
    Represents overall customer engagement and activity level
  \end{itemize}
\item
  \textbf{LD2} (4.2\% of variance) captures remaining separation

  \begin{itemize}
  \item
    Strong positive coefficient: average order value
  \item
    Strong negative coefficient: browsing time
  \item
    Represents order size patterns versus browsing efficiency
  \end{itemize}
\end{itemize}

The notebook's Module 4 displays both the discriminant coefficients
(scalings) and the group means on standardized features, allowing
interpretation of which behavioral patterns define each segment.

\textbf{Understanding Discriminant Coefficients}

Discriminant analysis produces two types of coefficients, each serving
distinct interpretive purposes:

\textbf{Unstandardized Coefficients}: These represent the raw weights
applied to predictor variables in their original measurement units. An
unstandardized coefficient of 0.5 for income (measured in thousands of
dollars) means that a 1,000 dollar increase in income increases the
discriminant score by 0.5 units. These coefficients depend on the
original measurement scales and cannot be directly compared across
variables measured in different units.

\textbf{Standardized Coefficients}: These coefficients indicate the
relative importance of each predictor independent of its original scale.
Standardization converts all predictors to have mean zero and standard
deviation one before computing discriminant functions. A standardized
coefficient of 0.8 for income and 0.3 for age suggests that income
contributes more strongly to group separation than age, after accounting
for differences in their variability. Standardized coefficients
facilitate comparison across predictors and are analogous to beta
coefficients in standardized regression.

The standardized form proves particularly valuable when comparing the
relative importance of predictors, while the unstandardized form is
necessary when applying the discriminant function to new observations in
their original measurement units.

\textbf{Discriminant Loadings}

Similar to factor analysis loadings - correlation between original
variable and discriminant function. Discriminant loadings represent the
simple correlation between each original predictor and the discriminant
scores, providing an alternative measure of variable importance.

High absolute loading = variable is important for that discriminant
function. Loadings above 0.30 in absolute value typically indicate
meaningful contributions to the discriminant function.

The relationship between coefficients and loadings mirrors factor
analysis: coefficients represent regression-like weights accounting for
correlations among predictors, while loadings represent simple
correlations ignoring other predictors. When predictors are highly
correlated, coefficients and loadings can differ substantially, with
loadings often providing more intuitive interpretation.

The notebook creates a DataFrame of coefficients to easily identify the
most influential features for each discriminant function.

\subsubsection{Step 5: Validation}

Comprehensive model validation ensures that discriminant analysis
results generalize beyond the training data and meet performance
requirements. Modules 3, 6, 9, and 10 of the accompanying notebook
provide detailed validation analyses that illustrate these principles.
Classification accuracy assessment begins with the confusion matrix
constructed from test set predictions, as presented in Module 9. This
matrix reveals not only overall accuracy but also class-specific
performance through precision, recall, and F1-scores compiled in the
classification report. Practitioners must resist the temptation to focus
solely on overall accuracy, as this aggregate metric can obscure poor
performance on specific groups, particularly when class sizes differ
substantially.

Cross-validation provides a more robust estimate of model performance by
evaluating the model across multiple train-test partitions of the data.
K-fold cross-validation, typically employing five or ten folds, assesses
performance stability across different data splits and guards against
the possibility that a single train-test split might yield
unrepresentative results. The notebook implementation utilizes the
\texttt{cross\_val\_score} function with five folds, generating a
distribution of accuracy estimates that characterizes model performance
more comprehensively than a single test set evaluation.

Visual validation techniques complement numerical performance metrics by
revealing the geometric structure of group separation. Module 5 presents
scatter plots of discriminant scores in the two-dimensional space
defined by LD1 and LD2, allowing visual assessment of how well the
discriminant functions separate groups. Module 8 provides decision
boundary plots that project the classification regions onto
two-dimensional feature subspaces, illustrating where the model assigns
observations to each group. Module 10 generates ROC curves for each
segment using the one-versus-rest approach, with the area under these
curves quantifying the model's ability to distinguish each segment from
all others. These visualizations verify that groups are well-separated
and reveal any regions of potential classification ambiguity.

Advanced validation metrics provide additional insights into model
performance characteristics. Module 10's ROC curve analysis with AUC
scores quantifies discrimination ability for each group independently,
revealing whether certain groups prove easier to classify than others.
Module 7 analyzes posterior probabilities to assess prediction
confidence, identifying observations where the model assigns high
probability to the predicted class versus those where probabilities
spread across multiple classes. Module 9's side-by-side confusion matrix
comparison for LDA versus QDA facilitates direct assessment of whether
QDA's additional complexity translates to meaningful performance
improvements.

\section{Applied Example: Marketing Segmentation}

\textbf{Complete Implementation}: See
\texttt{ch5\_guiding\_example/marketing\_discriminant\_analysis.ipynb}
for the full interactive analysis with visualizations and detailed
interpretations.

\subsection{Business Problem}

An e-commerce company has 1,200 customers and wants to classify them
into three segments for targeted marketing:

\begin{enumerate}
\item
  \textbf{High-Value} (30\%): Premium customers, high spending and
  engagement
\item
  \textbf{Loyal} (40\%): Regular customers, moderate spending,
  consistent
\item
  \textbf{Occasional} (30\%): Infrequent buyers, need re-engagement
\end{enumerate}

The dataset is synthetically generated using
\texttt{fetch\_marketing.py} to ensure reproducibility and known
statistical properties for educational purposes.

\subsection{Variables (p = 8)}

The analysis employs eight behavioral metrics that characterize customer
engagement patterns across multiple dimensions. Purchase frequency
measures the average number of transactions a customer completes per
month, providing insight into shopping regularity. Average order value,
denominated in US dollars, captures the typical monetary amount spent
per transaction, distinguishing high-spending customers from those
making smaller purchases. Browsing time quantifies the minutes spent per
website session, reflecting engagement level and product consideration
depth. Cart abandonment rate, expressed as a proportion between zero and
one, indicates the frequency with which customers initiate but fail to
complete purchases. Email open rate, similarly scaled from zero to one,
measures responsiveness to marketing communications. Loyalty points
represent the accumulated rewards a customer has earned through the
company's retention program. Support tickets count the average monthly
customer service interactions, potentially indicating product
satisfaction or purchase complexity. Social engagement tracks monthly
interactions with the company's social media presence, including likes,
shares, and comments.

The \texttt{MARKETING\_DATA\_DICTIONARY.md} document provides
comprehensive variable descriptions, including detailed range
specifications, measurement precision, and the multivariate normal data
generation methodology employed to create the synthetic dataset.

\subsection{Analysis Strategy}

\subsubsection{Why Both LDA and QDA?}

\textbf{LDA}: Assumes all customer segments have similar variability
patterns \textbf{QDA}: Allows different patterns (e.g., Occasional
customers might be more variable)

We'll fit both and compare performance.

\subsubsection{Feature Standardization}

Since variables are on different scales (dollars, rates, counts), we
standardize:

\[z_{j} = \frac{x_{j} - {\overline{x}}_{j}}{s_{j}}\]

This ensures no variable dominates due to scale.

\subsection{Results}

\textbf{Note}: The following results are from the Jupyter notebook
analysis. Run \texttt{marketing\_discriminant\_analysis.ipynb} to
reproduce these findings and generate visualizations.

\subsubsection{Discriminant Functions}

The notebook's Module 4 (Interpreting Discriminant Functions) reveals:

\textbf{LD1 (95.8\% of between-group variance)}:

\begin{itemize}
\item
  Separates High-Value from Occasional customers along the primary axis
\item
  Key drivers: Purchase frequency (strong negative), loyalty points
  (strong negative), average order value (negative)
\item
  Interpretation: Overall customer value and activity level
\end{itemize}

\textbf{LD2 (4.2\% of between-group variance)}:

\begin{itemize}
\item
  Distinguishes remaining group differences
\item
  Key drivers: Average order value (strong positive), browsing time
  (negative)
\item
  Interpretation: Order size versus browsing efficiency
\end{itemize}

\textbf{Insight}: Two independent dimensions describe customer segments:

\begin{enumerate}
\item
  Overall engagement and frequency (LD1 - dominant factor)
\item
  Purchase value patterns (LD2 - secondary factor)
\end{enumerate}

\subsubsection{Classification Performance}

\textbf{LDA Results} (Module 3):

\begin{itemize}
\item
  Test accuracy: Nearly perfect classification
\item
  Cross-validation: 99.9\% (± 0.3\%)
\item
  Strong performance across all segments
\end{itemize}

\textbf{QDA Results} (Module 6):

\begin{itemize}
\item
  Test accuracy: Perfect classification
\item
  Cross-validation: 100.0\% (± 0.0\%)
\item
  Slightly outperforms LDA with more flexible boundaries
\end{itemize}

\textbf{Interpretation}: Both models achieve excellent performance. The
synthetic data's well-separated structure allows for near-perfect
classification, demonstrating the power of discriminant analysis when
groups have distinct multivariate profiles.

\subsubsection{Model Comparison and Visualization}

The notebook includes comprehensive visualizations:

\textbf{Module 5}: Discriminant space scatter plot showing customer
distribution in LD1-LD2 space with group centroids

\textbf{Module 8}: QDA decision boundaries visualization using purchase
frequency and average order value

\textbf{Module 9}: Side-by-side confusion matrices comparing LDA and QDA
performance

\textbf{Module 10}: ROC curves for each segment showing excellent
discrimination (AUC near 1.0)

\subsubsection{Model Selection Recommendation}

Module 11 provides a comprehensive comparison and recommends
\textbf{LDA} despite QDA's marginally better performance because:

\begin{itemize}
\item
  Similar accuracy (difference is negligible)
\item
  Simpler model with fewer parameters
\item
  More interpretable discriminant functions
\item
  Lower overfitting risk
\item
  Easier to explain to stakeholders
\end{itemize}

\subsection{Business Insights}

\subsubsection{Segment Characteristics}

Module 4's group means analysis reveals distinct behavioral profiles
that characterize each customer segment. High-Value customers exhibit
positive standardized values across purchase frequency, order value, and
browsing time, indicating sustained engagement with the company's
offerings. This segment demonstrates exceptionally high email open rates
and social engagement, suggesting strong brand affinity and
responsiveness to marketing communications. Notably, these customers
show low cart abandonment rates and minimal support ticket generation,
reflecting purchase certainty and product satisfaction. The highest
loyalty points accumulation among this group confirms their long-term
value to the organization. The recommended strategy for High-Value
customers emphasizes retention through premium services, personalized
product recommendations, and exclusive benefits that reinforce their
valued status.

Loyal customers present moderate purchase frequency and order values,
positioning them between the extremes of occasional and high-value
segments. This group maintains good email engagement and accumulates
loyalty points steadily, though not at the accelerated pace observed
among High-Value customers. Their behavioral metrics demonstrate balance
across most dimensions, suggesting consistent but not exceptional
engagement. The strategic approach for Loyal customers focuses on
upselling opportunities that encourage larger purchases and
cross-selling initiatives that broaden their product adoption.
Enhancements to the loyalty program can incentivize increased spending
and engagement, potentially facilitating migration toward High-Value
status.

Occasional customers display negative standardized values across most
engagement metrics, indicating sporadic and limited interaction with the
company. This segment exhibits elevated cart abandonment rates and
generates more support tickets per capita, suggesting hesitation in
purchase decisions or difficulties in the buying process. Minimal
loyalty points accumulation and social engagement further characterize
this group's tenuous connection to the brand. The strategic imperative
for Occasional customers involves re-engagement campaigns designed to
increase purchase frequency, cart recovery mechanisms that address
abandonment triggers, and educational content that builds product
knowledge and purchase confidence.

\subsubsection{Actionable Applications}

The discriminant analysis framework enables several practical
applications that translate statistical insights into business value.
New customer classification becomes feasible once a customer accumulates
sufficient behavioral data, typically requiring two to three months of
transaction history. The model then automatically assigns the customer
to a segment, enabling tailored marketing strategies from an early stage
in the customer relationship. Module 7 demonstrates how posterior
probabilities provide confidence scores for each classification,
allowing the marketing team to gauge classification certainty and adjust
campaign intensity accordingly.

Monitoring segment migration over time reveals customer lifecycle
dynamics that inform retention strategies. Tracking discriminant scores
longitudinally identifies Occasional customers who transition toward
Loyal status, validating re-engagement efforts, as well as Loyal
customers whose scores drift toward Occasional levels, signaling
attrition risk. The discriminant scores serve as early warning
indicators of segment transition, enabling proactive interventions
before customers fully migrate to less valuable segments.

Marketing return on investment optimization leverages the classification
framework to allocate resources efficiently across customer segments.
Expensive retention campaigns targeting High-Value customers, where
classification confidence is typically high, maximize the impact of
premium marketing investments. Cost-effective email campaigns suit Loyal
customers, whose moderate value justifies ongoing engagement without
intensive personalization. Automated cart recovery mechanisms address
Occasional customers efficiently, deploying technology-enabled
interventions that require minimal manual effort.

Campaign personalization benefits from Module 7's posterior probability
analysis, which identifies customers with ambiguous segment membership
reflected in probability distributions spread across multiple classes.
These customers, exhibiting characteristics of multiple segments, may
respond favorably to hybrid marketing strategies that combine elements
designed for different segments, such as re-engagement messaging paired
with loyalty program incentives.

\subsection{Credit Risk Classification Example}

To illustrate discriminant analysis in a financial context, consider a
credit risk assessment problem where a bank classifies loan applicants
as either ``No Default'' (low risk) or ``Default'' (high risk) based on
financial characteristics.

\subsubsection{Problem Setup}

\textbf{Groups}: \(g = 2\)

\begin{itemize}
\item
  Group 0: No Default (historically 95\% of applicants)
\item
  Group 1: Default (historically 5\% of applicants)
\end{itemize}

\textbf{Predictors}: \(p = 3\)

\begin{itemize}
\item
  \(x_{1}\): Annual income (thousands of dollars)
\item
  \(x_{2}\): Debt-to-income ratio (proportion, 0 to 1)
\item
  \(x_{3}\): Credit score (300 to 850)
\end{itemize}

\subsubsection{Discriminant Score Interpretation}

After fitting LDA, suppose the discriminant function is:

\[\text{ Score } = 0.002 \cdot \text{ income } - 8.5 \cdot \text{ debt ratio } + 0.015 \cdot \text{ credit score } - 5.2\]

The cutoff threshold with equal prior probabilities would be zero.
However, recognizing the imbalanced class distribution and the higher
cost of default misclassification, the bank adjusts prior probabilities
to reflect business considerations rather than sample proportions.

\textbf{Example Classification}:

Consider applicant A with income equal to 50,000 dollars, debt ratio
equal to 0.40, and credit score equal to 650:

\[\text{ Score}_{A} = 0.002(50) - 8.5(0.40) + 0.015(650) - 5.2\]
\[\text{ Score}_{A} = 0.1 - 3.4 + 9.75 - 5.2 = 1.25\]

A positive discriminant score (1.25 greater than 0) indicates
classification into Group 0 (No Default), suggesting this applicant
presents low credit risk. The magnitude of the score reflects the
strength of classification confidence.

Consider applicant B with income equal to 35,000 dollars, debt ratio
equal to 0.65, and credit score equal to 580:

\[\text{ Score}_{B} = 0.002(35) - 8.5(0.65) + 0.015(580) - 5.2\]
\[\text{ Score}_{B} = 0.07 - 5.525 + 8.7 - 5.2 = - 1.955\]

A negative discriminant score (-1.955 less than 0) indicates
classification into Group 1 (Default), flagging this applicant as high
credit risk. The bank would likely deny this loan application or offer
modified terms with higher interest rates to compensate for elevated
risk.

\textbf{Coefficient Interpretation}:

\begin{itemize}
\item
  \textbf{Income} (+0.002): Higher income reduces default risk, though
  the small coefficient reflects income's measurement in thousands
\item
  \textbf{Debt Ratio} (-8.5): Strong negative coefficient indicates high
  debt loads substantially increase default risk
\item
  \textbf{Credit Score} (+0.015): Higher credit scores reduce default
  risk, reflecting proven creditworthiness
\end{itemize}

The debt-to-income ratio emerges as the strongest predictor when
coefficients are standardized, consistent with credit risk theory that
emphasizes debt burden as a primary default driver.

\textbf{Business Application}:

The discriminant score threshold can be adjusted to reflect business
priorities. Setting a higher threshold (e.g., requiring scores above 0.5
instead of 0) makes classification more conservative, approving fewer
borderline applicants but reducing default rates. Conversely, lowering
the threshold (e.g., -0.5) approves more applicants, increasing loan
volume but accepting higher default risk. Banks optimize this threshold
by analyzing the trade-off between foregone interest revenue from
rejected non-defaulters versus losses from accepted defaulters.

\section{Advanced Topics}

\subsection{Variable Selection}

The inclusion of all available predictors may not optimize discriminant
analysis performance, motivating variable selection procedures that
identify parsimonious models. Stepwise discriminant analysis employs
iterative algorithms to build or refine predictor sets. Forward
selection begins with an empty model and sequentially adds variables
that contribute most substantially to group separation, typically
measured by F-to-enter statistics or Wilks' Lambda changes. Conversely,
backward elimination starts with all candidate predictors and
iteratively removes those contributing least to discrimination, guided
by F-to-remove statistics. Both approaches utilize Wilks' Lambda or
related F-statistics to quantify how much each variable enhances group
separation. While computationally convenient, stepwise methods can be
unstable when predictors are highly correlated and may capitalize on
chance associations in finite samples.

Shrinkage methods offer modern alternatives that stabilize coefficient
estimates through regularization. Regularized Discriminant Analysis
(RDA) introduces a tuning parameter that shrinks the group-specific
covariance matrices toward a common structure, effectively interpolating
between LDA and QDA while reducing parameter variance. Penalized LDA
extends this concept by applying penalties such as the L1 norm to
discriminant coefficients, encouraging sparse solutions where many
coefficients equal zero. This sparsity proves valuable in
high-dimensional settings where interpretability benefits from
identifying a small subset of influential predictors.

\subsection{Handling Imbalanced Classes}

Substantial imbalance in group sizes, such as ninety-five percent
acceptable items versus five percent defective items, poses challenges
for discriminant analysis. When sample proportions serve as prior
probabilities, the classifier tends to favor the majority class,
achieving high overall accuracy by predominantly predicting the larger
group while performing poorly on the minority class. Two complementary
strategies address this issue.

Adjusting prior probabilities away from sample proportions recalibrates
the classification rule to account for the practical importance of
different groups. Setting equal priors regardless of sample sizes gives
each group equal a priori weight, preventing the majority class from
dominating predictions solely due to its prevalence. Alternatively,
priors can reflect business costs, assigning higher prior probability to
groups where misclassification carries greater consequences. For
instance, if failing to detect a defective item costs substantially more
than falsely flagging an acceptable item, elevated prior probability for
the defective class appropriately adjusts the decision threshold.

Sampling techniques modify the training data composition to mitigate
imbalance effects. Oversampling the minority class through replication
or synthetic data generation increases its representation in the
training set, allowing the model to learn its characteristics more
effectively. Undersampling the majority class reduces its dominance by
randomly selecting a subset of observations for model training. The
Synthetic Minority Over-sampling Technique (SMOTE) generates artificial
minority class observations by interpolating between existing minority
class instances in feature space, effectively expanding the minority
class representation while introducing variation beyond simple
replication.

\subsection{Model Diagnostics}

Several diagnostic statistics quantify the strength and significance of
group separation in discriminant analysis. Understanding these
diagnostics helps assess model quality and interpret the degree to which
groups differ in the multivariate space.

\subsubsection{Wilks' Lambda}

Wilks' Lambda tests whether group means differ significantly across the
predictor space, calculated as the ratio of the determinant of the
within-group sum of squares matrix \(\mathbf{W}\) to the determinant of
the total sum of squares matrix \(\mathbf{T}\):

\[\Lambda = \frac{|\mathbf{W}|}{|\mathbf{T}|}\]

Small values approaching zero indicate strong group separation, as the
within-group variation becomes negligible relative to total variation.
The test statistic ranges from 0 to 1, where:

\begin{itemize}
\item
  \(\Lambda \approx 0\): Perfect group separation (all variation is
  between groups)
\item
  \(\Lambda \approx 1\): No group separation (all variation is within
  groups)
\end{itemize}

The associated F-statistic provides formal hypothesis testing of whether
group means differ significantly. Wilks' Lambda is also used in stepwise
discriminant analysis as a criterion for variable selection, where
variables that minimize Lambda (maximize group separation) are preferred
for inclusion in the model.

\subsubsection{Hotelling's T-squared Test}

For testing whether two groups have different centroids (mean vectors),
Hotelling's T-squared test provides the multivariate generalization of
the univariate t-test. The test statistic is:

\[T^{2} = \frac{n_{1}n_{2}}{n_{1} + n_{2}}\left( \mathbf{\overline{x}}_{1} - \mathbf{\overline{x}}_{2} \right)^{\top}\mathbf{S}_{\text{pooled}}^{- 1}\left( \mathbf{\overline{x}}_{1} - \mathbf{\overline{x}}_{2} \right)\]

where \(n_{1}\) and \(n_{2}\) are the sample sizes for the two groups,
\(\mathbf{\overline{x}}_{1}\) and \(\mathbf{\overline{x}}_{2}\) are the
group mean vectors, and \(\mathbf{S}_{\text{pooled}}\) is the pooled
covariance matrix. This test determines whether the multivariate means
differ significantly, providing a foundation for discriminant analysis
by confirming that groups occupy distinct regions of the predictor
space.

For more than two groups, multivariate analysis of variance (MANOVA)
extends this concept, testing the null hypothesis that all group
centroids are equal. MANOVA and discriminant analysis are closely
related: MANOVA tests whether groups differ, while discriminant analysis
describes how they differ and builds classification rules based on those
differences.

\subsubsection{Eigenvalues and Discriminant Functions}

Each discriminant function has an associated eigenvalue that quantifies
its discriminatory power. The eigenvalue represents the ratio of
between-group variance to within-group variance along that discriminant
function:

\[\lambda_{i} = \frac{\text{between-group SS for LD}_{i}}{\text{ within-group SS for LD}_{i}}\]

Larger eigenvalues indicate stronger discrimination. A high eigenvalue
means that particular discriminant function effectively separates
groups, as the between-group variation dominates the within-group
variation along that dimension.

The proportion of discriminatory power explained by each function is
calculated as:

\[\text{ Proportion explained by LD}_{i} = \frac{\lambda_{i}}{\sum_{j = 1}^{m}\lambda_{j}}\]

where \(m = \min(g - 1,p)\) is the total number of discriminant
functions. Typically, the first discriminant function (LD1) accounts for
the largest proportion of group separation, with subsequent functions
capturing progressively less discrimination.

\textbf{Example Interpretation}: If LD1 has an eigenvalue of 12.5 and
LD2 has an eigenvalue of 0.8, then LD1 explains
\(\frac{12.5}{12.5 + 0.8} = 94\%\) of the discriminatory power. This
indicates that the first dimension captures nearly all meaningful group
separation, and LD2 contributes minimally.

\subsubsection{Canonical Correlation}

Canonical correlation measures the strength of the relationship between
discriminant functions and group membership. For each discriminant
function, the canonical correlation represents the square root of the
proportion of between-group variation to total variation:

\[R_{\text{can }} = \sqrt{\frac{\text{between-group SS}}{\text{ total SS}}} = \sqrt{\frac{\lambda_{i}}{1 + \lambda_{i}}}\]

Values approaching one indicate excellent discrimination, signifying
that the discriminant function effectively captures group differences.
Canonical correlation ranges from 0 to 1, where:

\begin{itemize}
\item
  \(R_{\text{can }} \approx 1\): Discriminant function perfectly
  separates groups
\item
  \(R_{\text{can }} \approx 0\): Discriminant function provides no group
  separation
\end{itemize}

Multiple discriminant functions yield multiple canonical correlations,
with the first function typically exhibiting the highest value. Squaring
the canonical correlation gives the proportion of variance in the
discriminant scores explained by group membership, analogous to
R-squared in regression.

\subsubsection{Mahalanobis Distance}

Mahalanobis distance quantifies the distance between an observation and
a group centroid, accounting for the covariance structure of the data.
For observation \(\mathbf{x}\) and group \(k\), the Mahalanobis distance
is:

\[D^{2}\left( \mathbf{x},\mathbf{\mu}_{k} \right) = \left( \mathbf{x} - \mathbf{\mu}_{k} \right)^{\top}\mathbf{\Sigma}^{- 1}\left( \mathbf{x} - \mathbf{\mu}_{k} \right)\]

This metric plays multiple roles in discriminant analysis:

\textbf{Classification}: With equal prior probabilities, LDA assigns
observations to the group with the smallest Mahalanobis distance to its
centroid. This geometric interpretation clarifies why LDA works: it
identifies the ``nearest'' group in multivariate space, where
``nearest'' accounts for variable correlations and scales.

\textbf{Outlier Detection}: Observations with large Mahalanobis
distances from all group centroids may be multivariate outliers
requiring investigation. Such observations lie far from typical group
members in the multivariate space.

\textbf{Variable Selection}: Mahalanobis distance between group
centroids can guide variable selection, as variables that increase the
distance between groups enhance discrimination.

Unlike Euclidean distance, Mahalanobis distance accounts for variable
correlations and unequal variances, making it more appropriate for
multivariate data where predictors are not independent and identically
scaled.

\subsection{Comparison with Other Methods}

Discriminant analysis occupies a particular niche among classification
methods, with its suitability depending on data characteristics and
analytical objectives. Logistic regression offers greater flexibility by
avoiding the normality assumption, modeling the log-odds of class
membership directly as a linear function of predictors. This approach
works particularly well for binary outcomes and provides probability
estimates through a natural probabilistic framework. However, extending
logistic regression to multiple groups through multinomial models
sacrifices some interpretability, as the method requires estimating
separate coefficient sets for each group contrast rather than generating
low-dimensional discriminant functions.

Support Vector Machines (SVM) provide powerful alternatives when
decision boundaries exhibit complex, non-linear structure. The kernel
trick enables SVMs to implicitly operate in high-dimensional feature
spaces where linear separation becomes feasible, accommodating intricate
decision boundaries without explicit feature engineering. SVMs impose no
distributional assumptions, making them robust across diverse data
types. The method's geometric focus on boundary optimization, however,
yields less interpretable models than discriminant functions, and SVMs
often excel in high-dimensional settings where discriminant analysis may
struggle due to parameter estimation challenges.

Random Forests handle non-linear relationships naturally through
recursive partitioning, building ensemble predictions from multiple
decision trees. This approach demonstrates robustness to outliers and
accommodates complex interactions without requiring explicit
specification. Variable importance measures derived from random forests
indicate which predictors contribute most to classification accuracy.
The ensemble nature of random forests, however, creates black-box models
that resist interpretation beyond variable importance rankings, limiting
their utility when understanding group differences constitutes a primary
analytical goal.

Discriminant analysis proves most appropriate when several conditions
align. Moderate sample sizes and dimensionality suit the method's
parameter estimation requirements, as extreme dimensionality or limited
sample sizes introduce estimation instability. Interpretability
requirements favor discriminant analysis, as the discriminant functions
directly reveal which variable combinations separate groups. Analytical
objectives focused on understanding group differences rather than
maximizing predictive accuracy align with the method's strengths.
Finally, data conforming reasonably well to multivariate normality with
similar within-group covariance structures satisfy the method's
parametric assumptions, ensuring optimal performance.

\section{Common Pitfalls and Best Practices}

\subsection{Common Mistakes}

Several pitfalls frequently undermine discriminant analysis
applications, warranting careful attention throughout the modeling
process. Ignoring fundamental assumptions represents a primary source of
error. Applying LDA when groups exhibit clearly different covariance
structures violates the method's homogeneity assumption, potentially
yielding biased decision boundaries that perform poorly on new data.
Outliers exert disproportionate influence on discriminant functions due
to their impact on estimated means and covariances, yet analysts
sometimes proceed without outlier screening, particularly when sample
sizes appear large. Mahalanobis distance calculations identify
multivariate outliers that may not be apparent from univariate
examinations.

The normality assumption requires that predictors follow multivariate
normal distributions within each group. While discriminant analysis
demonstrates robustness to moderate violations when sample sizes are
large, severe departures from normality can compromise both
classification accuracy and the optimality properties of the Bayes rule.
Particularly problematic are categorical predictors, which inherently
violate the continuity assumption underlying multivariate normality.
Including categorical variables (such as gender coded as 0/1, or product
category coded as 1/2/3) in discriminant analysis creates theoretical
inconsistencies, as discrete distributions cannot be multivariate
normal. When categorical predictors are essential for classification,
analysts face several alternatives: convert categories to quantitative
proxies when meaningful (e.g., replace product category with average
price), use logistic regression which accommodates categorical
predictors naturally, or employ non-parametric classification methods
such as classification trees or k-nearest neighbors that impose no
distributional assumptions.

Overfitting poses particular risks when the number of predictors
approaches or exceeds the sample size relative to the number of groups.
The discriminant analysis literature suggests requiring at least twenty
observations per predictor per group to ensure stable parameter
estimates. Violating this guideline risks producing discriminant
functions that fit idiosyncrasies of the training sample rather than
capturing genuine group differences, leading to poor generalization
performance.

The practice of evaluating model performance on training data yields
overly optimistic accuracy estimates that mislead analysts regarding
true predictive performance. Training accuracy necessarily exceeds test
accuracy because the model parameters are optimized to fit the training
observations. Proper validation requires held-out test sets or
cross-validation procedures that assess performance on observations not
used for parameter estimation.

Class imbalance creates a subtle trap where models achieve high overall
accuracy by predominantly predicting the majority class while failing to
identify minority class members effectively. Overall accuracy masks this
problem, as correct classification of the majority class dominates the
aggregate metric. Examining per-class precision, recall, and F1-scores
reveals whether the model performs adequately across all groups rather
than succeeding primarily through majority class predictions.

Correlated predictors introduce multicollinearity that inflates
coefficient standard errors and produces unstable discriminant functions
where small data perturbations yield substantially different
coefficients. The presence of highly redundant variables provides no
additional information for discrimination while complicating
interpretation and increasing estimation variance. Variable selection or
principal component analysis can address severe multicollinearity by
reducing predictor dimensionality.

\subsection{Best Practices}

Successful discriminant analysis implementation rests on several
foundational practices that enhance reliability and interpretability.
Data quality requires meticulous attention before model fitting begins.
Missing data must be handled appropriately, either through principled
imputation methods or case deletion when missingness is minimal and
random. Outlier screening using Mahalanobis distance identifies
observations that deviate substantially from their group's multivariate
center, allowing assessment of whether these points represent data entry
errors, measurement anomalies, or genuine extreme cases warranting
special handling. Verification of data entry errors prevents spurious
results arising from typographical mistakes or unit conversion errors.

Model selection proceeds most effectively through a disciplined
progression from simplicity to complexity. Beginning with LDA
establishes a baseline that subsequent, more complex models must
meaningfully improve upon to justify their adoption. Cross-validation
provides honest performance estimates that account for sampling
variability, generating multiple train-test partitions to characterize
performance distributions rather than relying on a single partition that
might be unrepresentative. Multiple performance metrics paint a
comprehensive picture of model behavior, with accuracy, precision,
recall, F1-scores, and AUC each illuminating different aspects of
classification performance.

Interpretation extends beyond merely reporting classification accuracy
to explaining the substantive meaning of discriminant functions.
Coefficient examination reveals which predictors drive group separation
and how they combine to form discriminant scores. Visualizing decision
boundaries, when feasible through dimensionality reduction, provides
geometric intuition about classification regions. Translating
statistical results into domain-specific insights ensures that
stakeholders understand not just which observations are classified into
which groups, but why the classification occurs and what it means for
practical decision-making.

Validation practices maintain model reliability over time and across
contexts. Testing on truly independent data, distinct from any data used
during model development, provides the most conservative performance
assessment. Monitoring performance over time in production deployments
detects concept drift, where the relationships between predictors and
groups evolve, degrading model accuracy. Updating models as patterns
change ensures continued relevance, with retraining schedules determined
by the rate of observed performance degradation.

\section{Summary and Key Takeaways}

Discriminant analysis provides a classification framework for assigning
observations to predefined groups based on measured predictor variables.
The method identifies linear combinations of predictors in LDA or
quadratic combinations in QDA that maximally separate groups in the
feature space. Under the assumptions of multivariate normality and known
covariance structures, discriminant analysis implements the optimal
Bayes classification rule, minimizing the total probability of
misclassification when these conditions hold.

The method proves most appropriate when analysts possess labeled
training data with known group memberships, seek interpretable insights
into group differences rather than black-box predictions, work with
moderate sample sizes and dimensionality where parameter estimation
remains stable, and observe data that conform reasonably well to
multivariate normality assumptions. These conditions position
discriminant analysis as a valuable tool when understanding group
structure constitutes a primary analytical objective alongside
prediction.

Several key decisions shape discriminant analysis applications. The
choice between LDA and QDA balances the parsimony of equal covariance
assumptions against the flexibility of group-specific covariances, with
LDA serving as the recommended starting point. Predictor selection
combines domain knowledge about theoretically relevant variables with
statistical criteria that identify redundant or non-discriminating
features. Prior probability specification chooses between equal priors
that weight groups uniformly, proportional priors that reflect sample or
population frequencies, or cost-based priors that account for asymmetric
misclassification consequences.

Interpretation distinguishes discriminant analysis from purely
predictive methods. The discriminant functions reveal how groups differ
by showing which variable combinations drive separation. Discriminant
loadings or coefficients indicate which variables contribute most
substantially to each function. Classification accuracy quantifies
predictive performance, but the interpretive insights into group
structure often provide greater analytical value than prediction
accuracy alone.

Validation requirements apply universally to discriminant analysis
applications. Hold-out test sets or cross-validation procedures provide
honest performance estimates uncontaminated by training data
optimization. Per-group performance metrics ensure that classification
succeeds across all groups rather than primarily through majority class
prediction. Monitoring performance over time in operational deployments
detects degradation and informs model updating schedules.

The power of discriminant analysis ultimately lies not merely in
classification accuracy, but in understanding what makes groups
different. This interpretive capability transforms multivariate data
into actionable insights that inform both scientific understanding and
practical decision-making, distinguishing the method from alternative
classification approaches that prioritize prediction over explanation.

\section{Putting It Into Practice}

To solidify your understanding of discriminant analysis, work through
the complete marketing segmentation example:

\textbf{Step 1}: Generate the customer dataset

\begin{Shaded}
\begin{Highlighting}[]
\BuiltInTok{cd}\NormalTok{ ch5\_guiding\_example}
\ExtensionTok{python}\NormalTok{ fetch\_marketing.py}
\end{Highlighting}
\end{Shaded}

\textbf{Step 2}: Open and run the Jupyter notebook

\begin{Shaded}
\begin{Highlighting}[]
\ExtensionTok{jupyter}\NormalTok{ notebook marketing\_discriminant\_analysis.ipynb}
\end{Highlighting}
\end{Shaded}

\textbf{Step 3}: Execute each module sequentially, reading the
explanations and examining the outputs

\textbf{Step 4}: Experiment with modifications:

\begin{itemize}
\item
  Try different train/test splits
\item
  Exclude certain features to see impact on performance
\item
  Adjust prior probabilities in the discriminant models
\item
  Create additional visualizations
\end{itemize}

\textbf{Step 5}: Review the generated PNG files:

\begin{itemize}
\item
  \texttt{marketing\_lda\_scores.png}: Discriminant space visualization
\item
  \texttt{marketing\_qda\_boundaries.png}: Decision boundaries in 2D
\item
  \texttt{marketing\_confusion\_comparison.png}: LDA vs QDA performance
\item
  \texttt{marketing\_roc\_curves.png}: Segment-specific classification
  quality
\end{itemize}

By working through the complete pipeline from data generation to model
comparison, you will develop practical skills in applying discriminant
analysis to real-world classification problems.

\section{AI-Assisted Content Development}

These companion notes were developed with the assistance of artificial
intelligence (Claude, Anthropic) to enhance their pedagogical value and
comprehensiveness. The AI contribution involved:

\textbf{Content Organization and Clarity}: Restructuring technical
concepts into a logical progression from foundational theory to advanced
applications, ensuring accessibility for students with diverse
mathematical backgrounds.

\textbf{Mathematical Exposition}: Translating formal statistical
notation into clear explanations with concrete interpretations,
balancing mathematical rigor with practical understanding.

\textbf{Example Development}: Creating worked examples (marketing
segmentation, credit risk classification) that demonstrate discriminant
analysis application to realistic business problems, complete with
numerical calculations and business interpretations.

\textbf{Coverage Alignment}: Enhancing content to ensure comprehensive
coverage of assessment topics, including detailed discussions of
diagnostic statistics (Wilks' Lambda, Hotelling's T-squared,
eigenvalues, canonical correlation, Mahalanobis distance), coefficient
interpretation (standardized vs unstandardized), and methodological
considerations (categorical predictors, assumption violations).

\textbf{Technical Accuracy}: Verifying formulas, statistical concepts,
and Python implementation details against authoritative references to
ensure correctness.

The pedagogical structure, learning objectives, and example selection
reflect the course instructor's vision (Dr. Juliho Castillo), with AI
serving as a tool for content development and refinement rather than as
the primary author. All technical content has been reviewed for accuracy
and alignment with course objectives.

\textbf{AI Tool Used}: Claude 3.5 Sonnet (Anthropic) \textbf{Development
Dates}: October 2024 - October 2025 \textbf{Human Oversight}: Dr. Juliho
Castillo, Tecnologico de Monterrey

This disclosure reflects the growing role of AI in educational content
creation while maintaining transparency about the development process.
Students are encouraged to engage critically with all course materials,
verify understanding through practice problems, and seek clarification
from the instructor when concepts require additional explanation.

{ ``The goal is to turn data into information, and information into
insight.''\\
- Carly Fiorina }

\end{document}

